\chapter{System Workflows and Data Flow}

\section{Authentication Workflow}

The authentication process is a critical component of the Campus-Sync platform, ensuring secure access while providing a seamless user experience.

\subsection{User Login Process}

\begin{enumerate}
    \item User navigates to the authentication page
    \item User enters credentials (email and password)
    \item Frontend validates input format
    \item Credentials are sent to the backend authentication API
    \item Backend validates credentials against the database
    \item Upon successful validation, JWT tokens are generated
    \item Tokens are returned to the frontend
    \item Frontend stores tokens in secure storage
    \item User is redirected to their role-specific dashboard
\end{enumerate}

\begin{figure}[H]
    \centering
    \begin{tikzpicture}[thick, node distance=2.5cm, every node/.style={align=center, font=\small}]
        \node (user) [draw, rectangle, minimum width=2.5cm, minimum height=1cm] {User};
        \node (frontend) [draw, rectangle, minimum width=2.5cm, minimum height=1cm, right=of user] {Frontend};
        \node (backend) [draw, rectangle, minimum width=2.5cm, minimum height=1cm, right=of frontend] {Backend};
        \node (database) [draw, rectangle, minimum width=2.5cm, minimum height=1cm, right=of backend] {Database};
        
        \draw[->] (user) -- (frontend) node[midway, above, font=\scriptsize, yshift=0.3cm] {1. Enter\\credentials};
        \draw[->] (frontend) -- (backend) node[midway, above, font=\scriptsize, yshift=0.3cm] {2. Submit auth\\request};
        \draw[->] (backend) -- (database) node[midway, above, font=\scriptsize, yshift=0.3cm] {3. Validate\\credentials};
        \draw[->] (database) -- (backend) node[midway, above, font=\scriptsize, yshift=0.3cm] {4. User\\data};
        \draw[->] (backend) -- (frontend) node[midway, above, font=\scriptsize, yshift=0.3cm] {5. JWT\\tokens};
        \draw[->] (frontend) -- (user) node[midway, above, font=\scriptsize, yshift=0.3cm] {6. Redirect to\\dashboard};
    \end{tikzpicture}
    \caption{Authentication Flow}
    \label{fig:auth-flow}
\end{figure}

\subsection{Token Management}

\begin{itemize}
    \item Access tokens with 15-minute expiry for security
    \item Refresh tokens with 7-day expiry for convenience
    \item Automatic token refresh before expiry
    \item Secure storage of tokens in HTTP-only cookies
\end{itemize}

\section{Assignment Management Workflow}

The assignment management system facilitates the creation, distribution, and evaluation of academic assignments.

\subsection{Assignment Creation Process}

\begin{enumerate}
    \item Teacher navigates to assignment creation interface
    \item Teacher fills in assignment details (title, description, due date, etc.)
    \item Teacher attaches any necessary files or resources
    \item Teacher selects target course and students
    \item Assignment is saved to the database
    \item Notifications are sent to relevant students
    \item Assignment appears in students' assignment lists
\end{enumerate}

\subsection{Assignment Submission Process}

\begin{enumerate}
    \item Student views assignment in their assignment list
    \item Student prepares submission (writes response, attaches files)
    \item Student submits assignment through the platform
    \item Submission is stored in the database
    \item Teacher receives notification of new submission
    \item Submission appears in teacher's evaluation queue
\end{enumerate}

\begin{figure}[H]
    \centering
    \begin{tikzpicture}[thick, node distance=2.5cm, every node/.style={align=center, font=\small}]
        \node (student) [draw, rectangle, minimum width=2.5cm, minimum height=1cm] {Student};
        \node (frontend) [draw, rectangle, minimum width=2.5cm, minimum height=1cm, right=of student] {Frontend};
        \node (backend) [draw, rectangle, minimum width=2.5cm, minimum height=1cm, right=of frontend] {Backend};
        \node (database) [draw, rectangle, minimum width=2.5cm, minimum height=1cm, right=of backend] {Database};
        \node (teacher) [draw, rectangle, minimum width=2.5cm, minimum height=1cm, below=of frontend, yshift=-1.5cm] {Teacher};
        
        \draw[->] (student) -- (frontend) node[midway, above, font=\scriptsize, yshift=0.3cm] {1. Submit\\assignment};
        \draw[->] (frontend) -- (backend) node[midway, above, font=\scriptsize, yshift=0.3cm] {2. Upload file\\+ metadata};
        \draw[->] (backend) -- (database) node[midway, above, font=\scriptsize, yshift=0.3cm] {3. Store\\assignment data};
        \draw[->] (database) -- (backend) node[midway, above, font=\scriptsize, yshift=0.3cm] {4.\\Confirmation};
        \draw[->] (backend) -- (frontend) node[midway, above, font=\scriptsize, yshift=0.3cm] {5. Success\\response};
        \draw[->] (frontend) -- (student) node[midway, above, font=\scriptsize, yshift=0.3cm] {6. Show\\confirmation};
        \draw[->] (backend) -- (teacher) node[midway, right, font=\scriptsize, xshift=0.3cm] {7. Notification\\(async)};
    \end{tikzpicture}
    \caption{Assignment Submission Flow}
    \label{fig:assignment-flow}
\end{figure}

\section{Grade Management Workflow}

The grade management system provides teachers with tools to evaluate student performance and students with access to their academic progress.

\subsection{Grade Entry Process}

\begin{enumerate}
    \item Teacher accesses grade entry interface
    \item Teacher selects course and assignment/exam
    \item Teacher enters grades for individual students
    \item System validates grade format and range
    \item Grades are saved to the database
    \item Grade calculations are updated (averages, totals)
    \item Students receive notifications of grade updates
\end{enumerate}

\subsection{Grade Viewing Process}

\begin{enumerate}
    \item Student navigates to grade viewing section
    \item System retrieves student's grades from database
    \item Grades are organized by course and term
    \item Visual representations (charts, graphs) are generated
    \item Student views grades and performance metrics
\end{enumerate}

\section{Attendance Tracking Workflow}

The attendance system helps track student participation and engagement in classes.

\subsection{Attendance Recording Process}

\begin{enumerate}
    \item Teacher opens attendance interface for a specific class
    \item System displays enrolled students for that class
    \item Teacher marks attendance status for each student
    \item Attendance data is saved to the database
    \item Attendance statistics are updated
    \item Notifications may be sent for attendance concerns
\end{enumerate}

\subsection{Attendance Viewing Process}

\begin{enumerate}
    \item Student or admin accesses attendance records
    \item System retrieves attendance data from database
    \item Data is organized by course and date
    \item Attendance percentages and trends are calculated
    \item Visual reports are generated
\end{enumerate}

\section{Billing and Payment Workflow}

The financial management system handles tuition, fees, and other monetary transactions.

\subsection{Billing Process}

\begin{enumerate}
    \item Admin creates billing records for students
    \item System calculates amounts based on course enrollment
    \item Billing records are stored in the database
    \item Invoices are generated and sent to students
    \item Payment due dates are tracked
\end{enumerate}

\subsection{Payment Process}

\begin{enumerate}
    \item Student accesses billing and payment interface
    \item Student views outstanding bills and due dates
    \item Student selects payment method and amount
    \item Payment information is processed securely
    \item Payment confirmation is recorded in the database
    \item Receipts are generated and sent to student
\end{enumerate}

\section{Communication Workflow}

The communication system facilitates announcements, notifications, and messaging between users.

\subsection{Announcement Creation Process}

\begin{enumerate}
    \item Admin or teacher creates an announcement
    \item Announcement details are entered (title, content, audience)
    \item Target audience is selected (all users, specific roles, etc.)
    \item Announcement is saved to the database
    \item Notifications are sent to target audience
    \item Announcement appears in relevant users' feeds
\end{enumerate}

\subsection{Notification Delivery Process}

\begin{enumerate}
    \item System event triggers a notification
    \item Notification is created and stored in database
    \item Notification is formatted for delivery
    \item Notification is sent via preferred channels (in-app, email, SMS)
    \item User receives notification
    \item User interaction with notification is tracked
\end{enumerate}

\section{Data Synchronization}

Campus-Sync implements real-time data synchronization to ensure all users have access to the most current information.

\subsection{Real-Time Updates}

\begin{itemize}
    \item WebSocket connections for live data updates
    \item Database triggers for immediate change notifications
    \item Client-side caching for improved performance
    \item Conflict resolution for concurrent updates
\end{itemize}

\subsection{Offline Support}

\begin{itemize}
    \item Local data storage for offline access
    \item Synchronization queue for offline actions
    \item Conflict detection and resolution
    \item Automatic sync when connectivity is restored
\end{itemize}

\section{Error Handling and Recovery}

The system implements comprehensive error handling to ensure reliability and user experience.

\subsection{Error Types}

\begin{itemize}
    \item Validation errors (input format issues)
    \item Authentication errors (login failures)
    \item Authorization errors (access denied)
    \item Database errors (connection failures, query issues)
    \item Server errors (500-level HTTP errors)
    \item Network errors (connectivity issues)
\end{itemize}

\subsection{Error Response Format}

\begin{lstlisting}[language=JSON, caption=Error Response Format]
{
  "success": false,
  "error": {
    "code": "VALIDATION_ERROR",
    "message": "Validation failed",
    "details": [
      {
        "field": "email",
        "message": "Email is required"
      }
    ]
  }
}
\end{lstlisting}

\subsection{Recovery Mechanisms}

\begin{itemize}
    \item Automatic retry for transient failures
    \item Graceful degradation for non-critical features
    \item User-friendly error messages
    \item Logging for debugging and monitoring
\end{itemize}

This workflow documentation provides a comprehensive overview of how data flows through the Campus-Sync platform and how various processes are executed, ensuring a clear understanding of system operations for developers, administrators, and stakeholders.